\subsection{Definition-driven normalization}
\label{sec:normalization-algorithm}

Normalization changes the representation of the current state, not the set of
executions being tested. Its input is the assertion $P$, the tuple $\bar t$ of
terms used by the continuation, a definition environment $\Gamma$, and a finite
budget $b$. The environment records qualified definitions, constructor and
binder types, and the interpretation of spatial notation. An admitted rewrite
must retain the joint observations of $P$ and $\bar t$.

\paragraph{Admission and branch selection.}
For a definition $H(\bar x)$, let a clause supply a guard $g_i$, fresh witnesses
$\bar w_i$, and body $B_i$. A licensed clause establishes, in the selected
interpretation,
\[
 g_i(\bar a) \Longrightarrow
 \bigl(H(\bar a)\ \equiv\ \exists\bar w_i.\,B_i[\bar a/\bar x]\bigr).
\]
The body includes every pure conjunct, spatial conjunct and typed existential;
a heap-only consequence of the body is insufficient. Constructor matching
can establish the guard without solving a new logical problem. Otherwise a
branch is selected only when its guard is established from the current state,
using admitted computation or checked reasoning. Pending annotations cannot
supply branch-selection assumptions. Unknown guards retain the application;
a symbolic case split is sound only if its alternatives cover the original
state and all branch conditions remain explicit.

Quoted logical conjunctions are parsed by nesting depth, outside field-name
strings. A supported $G_1\land G_2$ is lowered to separate pure conjuncts,
using $pure(G_1\land G_2)\equiv pure(G_1)\land pure(G_2)$ in the selected
assertion interpretation. All conjuncts and typed witnesses remain in the
clause; an unsupported conjunct rejects the clause rather than disappearing.
The generated full branch equality checks this lowering against the original
definition, including any library-specific notation interpretation.

Closed arithmetic and Boolean operations use their defined interpretation.
A known constructor enables match reduction, and a nonrecursive definition can
be unfolded with symbolic arguments if its complete body is representable.
A recursive definition unfolds when its selected branch is available; known
list or tree constructors often expose a smaller recursive argument. A
symbolic pointer does not prevent unfolding when the logical shape is known.
Neither heap ownership nor an unproved strategy consequence licenses a new
pure fact. Unknown definitions, unsupported binders or operators, and unresolved
layout interpretations remain residual terms.

\paragraph{Shared-base offsets and list slices.}
A symbolic base does not prevent evaluating a difference of known offsets.
For structurally identical residual bases, normalization applies
\[
 (b+u)-(b+v)=u-v,\qquad
 \operatorname{quot}((b+is)-(b+js),s)=i-j\quad(s\ne0).
\]
Offsets are computed with arbitrary-precision constants; a result outside the
AST constant representation is retained. Different residual bases are not
cancelled. These logical equalities retain the executor's C safety obligations;
they do not establish that unrelated pointers may be subtracted.

The standard QCP list primitive uses
$sublist(lo,hi,L)=skipn(Z.to\_nat(lo),firstn(Z.to\_nat(hi),L))$.
For a known constructor spine of length $m$, set
$a=\min(m,\max(0,lo))$ and $z=\min(m,\max(0,hi))$.
Evaluation returns the original elements at indices $a,\ldots,z-1$, or the
empty list when $z\le a$. Elements may remain symbolic; constructors, element
types and ordering are preserved. Unknown spines remain applications. The
implementation limits spine inspection to 4096 nodes and clamps large bounds
before conversion to a machine integer. In particular, negative lower bounds
follow QCP's two conversions rather than a naive length $hi-lo$.

\begin{lstlisting}[language={},numbers=left,caption={Budgeted normalization contract.}]
Normalize(Gamma, P, liveTerms, budget):
  while budget > 0:
    choose an occurrence with an admitted reduction
    if no occurrence exists: return (P, liveTerms, residuals)
    if the reduction unfolds a definition:
      resolve the original qualified definition
      match evaluated arguments; establish the clause guard
      freshen all local binders, preserving their types
      substitute arguments into the complete clause body
    replace the occurrence by its licensed equivalent
    rewrite every dependent live term consistently
    simplify newly exposed admitted computations
    budget := budget - 1
  return (P, liveTerms, remaining occurrences)
\end{lstlisting}

Each admitted rewrite consumes budget, including recursive unfolding. Thus the
algorithm terminates even when structural decrease cannot be established. An
unexpanded application retains its original meaning, so budget exhaustion
reduces automation rather than establishing a successful test. Rewriting an
unused existential requires a witness of its type; the algorithm cannot
silently discard a possibly empty nominal domain. The implementation's
constructor patterns and expression evaluator instantiate this contract for
their supported fragment. Definition and interpretation certificates license
external clauses; the parser and strategy encoding remain implementation
correspondence obligations. The algorithm does not certify arbitrary imported
strategies as equivalences.

A complementary reduction exposes a checked pure consequence without
replacing the spatial predicate. For a clause $H(\bar a)$ and a pure condition
$G(\bar a)$ independent of its local existential witnesses, the translator
proposes
\[
 H(\bar a) \equiv H(\bar a)\land G(\bar a).
\]
The kernel checks this equality against the original definition. Retaining
$H$ preserves its resources and witnesses; no heap-layout conversion is needed
when $G$ contains no field addresses. The generated rule adds the condition
only when absent, both at normalization points and between proof-search steps
that expose further recursive applications. Pure-fact closure delays removal of tautologies until the
closure finishes, preventing repeated insertion of a condition such as
$4096\ne0$. This mechanism can preserve facts before an imported proof-search
strategy weakens an assertion. It does not certify that subsequent strategy as
an equivalence or establish the completeness of predicate unfolding.

For example, a nonempty list clause expands to
\[
 sll(p,a::xs)\equiv
 p\ne0\ \land\ \exists q.\;
 data(p,a)*next(p,q)*sll(q,xs).
\]
Removing $p\ne0$ yields only a weaker assertion and can make a subsequent
folding VC unprovable. The same requirement applies to bounds, ordering
relations, tree invariants and abstract pure predicates in any definition;
it is independent of the predicate's name.

For a logical environment $\eta$, state $s$ and observed tuple $\bar v$, define
\[
 \operatorname{Obs}(P,\bar t)(s,\bar v)
 = \exists\eta.\;P(s,\eta)\land
             \bar v=\operatorname{eval}(\bar t)_{s,\eta}.
\]
The normalization obligation is equality of this relation before and after
rewriting. Composition of licensed steps preserves the relation. Consequently
it preserves initial nonemptiness and the validity of obligations over the
represented observations. These are conditional semantic results: connecting
actual AST rewrites, layouts and substitutions to the licensed steps is a
separate implementation task. A residual obligation is sent unchanged to QCP's
solvers, and to independently checked ITP when automatic reasoning remains
inconclusive.
